\subsection{A four-local Gibbs
realization}\label{proof:thin_shell_obstruction:a-four-local-gibbs-realization}

Set

\[
H_n(x)=n(S(x)^2-16)^2,\qquad
\widetilde\mu_n(x)=\frac{e^{-H_n(x)}}{\sum_y e^{-H_n(y)}}.
\tag{4.1}
\]

The inverse temperature is \(\beta =1\). The ground configurations are
exactly \(S=\pm 4\), and every finite-n Gibbs weight is positive. Since
n is even, off these shells the least value of \((S^{2}-16)^{2}\) is
144, attained at \(S=\pm 2\). Every off-shell configuration has energy
at least \(144n\).

Writing \(\delta_n=\widetilde\mu_n(S\notin\{-4,4\})\), the zero-energy
configurations give a partition-function lower bound
\(2\binom{n}{n/2+2}\). Hence

\[
\delta_n\le\frac{2^n e^{-144n}}{2\binom n{n/2+2}}
\le2^{n-1}e^{-144n}\le e^{-143n}.
\tag{4.2}
\]

Conditional on the ground shells the law is \(\nu _n\). Symmetry and
(1.3) therefore yield

\[
\operatorname{Var}_{\widetilde\mu_n}(f_n)\ge(1-\delta_n)n^{-10},
\qquad \max_i I_i^{\widetilde\mu_n}(f_n)\le\delta_n\le e^{-143n}.
\tag{4.3}
\]

\textbf{Theorem T2.} The fixed-temperature Gibbs family (4.1), generated
by interactions involving at most four spins, fails every polynomial AA
influence bound.

The proof is the exponential-versus-polynomial comparison in T1. This is
an actual low-arity Hamiltonian counterexample, rather than failure of
an available proof technique.

\subsubsection{Exact expansion and interaction
strength}\label{proof:thin_shell_obstruction:exact-expansion-and-interaction-strength}

On the Boolean cube,

\[
S^4=3n^2-2n+(12n-16)\sum_{i<j}x_ix_j
+24\sum_{i<j<k<\ell}x_ix_jx_kx_\ell.
\]

It follows that

\[
H_n(x)=n(3n^2-34n+256)
+n(12n-80)\sum_{i<j}x_ix_j
+24n\sum_{i<j<k<\ell}x_ix_jx_kx_\ell.
\tag{4.4}
\]

The constant does not affect the law. The nonconstant interactions are
two-local and four-local, but their strengths grow with n.~The maximum
incident absolute interaction strength equals

\[
\Lambda_n=(n-1)|n(12n-80)|+24n\binom{n-1}{3}=\Theta(n^4).
\]

The energy range is \(n(n^{2}-16)^{2}=\Theta (n^{5})\) for even
\(n\ge 6\). \textbf{This is not an extensive Hamiltonian with bounded
local strength and does not satisfy the row/local-field normalization of
the bounded-degree Ising examples.} Rescaling changes effective
temperature; the conclusion cannot be transferred to that different
normalization. The example disproves criteria based on interaction
arity, covariance norm, or magnetization concentration alone.

\subsubsection{Covariance remains bounded in the Gibbs
realization}\label{proof:thin_shell_obstruction:covariance-remains-bounded-in-the-gibbs-realization}

The elementary bounds
\(16(1-\delta _n)\le \mathbb{E} S^{2}\le 16+n^{2}\delta _n\) already
give \(\operatorname{Var}(M)\sim 16/n^{2}\) and a uniform covariance
bound. In fact \(\mathbb{E} S^{2}<16\) for every even \(n\ge 6\), so the
same norm bound two holds for \(n\ge 8\).

To verify this, examine the unnormalized sum of \((S^{2}-16)e^{-H_n}\).
The two shells \(S=\pm 2\) contribute at most \(-24e^{-144n}\), even
ignoring their binomial multiplicities. Positive contributions come only
from \(|S|\ge 6\), where \(H_n\ge 400n\), and their total is at most
\(n^{2}2^n e^{-400n}\), which is strictly smaller than \(24e^{-144n}\).
The \(S=0\) contribution is negative and \(S=\pm 4\) contributes zero.
Therefore the sum is negative.

For any exchangeable mean-zero spin law, the covariance eigenvalues are
\(\mathbb{E} S^{2}/n\) and \((n^{2}-\mathbb{E} S^{2})/(n(n-1))\). Here
they are at most \(16/n\) and \(n/(n-1)\), respectively, proving
\(\lVert \operatorname{Cov}(X)\rVert_{\mathrm{op}}\le 2\) for even
\(n\ge 8\).
